\section{Space Complexity --- Foundations}
%%% generated by the blueprint pipeline from TCSlib.Complexity.SpaceComplexity.Basic
\subsection*{Space Complexity --- Foundations: Introduction}

Declarations of the Lean module \texttt{TCSlib.Complexity.SpaceComplexity.Basic}.

\subsection*{Space Complexity --- Foundations: Results}

\begin{definition}[Deterministic space complexity class]
\lean{Complexity.SPACE}
\statementsource{Sav70}{
  pdf-c8033d2fd705-p002-b002,
  pdf-c8033d2fd705-p003-b001,
  pdf-c8033d2fd705-p003-b002
}
\statementsource{BGW70}{
  pdf-b9fb35a232af-p003-b001,
  pdf-b9fb35a232af-p004-b001,
  pdf-b9fb35a232af-p005-b001
}
\statementsource{BGS75}{pdf-2459e3840724-p004-b006}
\label{Complexity.SPACE}
\leanok
\uses{Turing.FinTM, Turing.FinTM.DecidesInSpace}
For a function $s:\mathbb{N}\to\mathbb{N}$, let $\mathrm{SPACE}(s)$ be the class of
Boolean languages $L$ for which there exist a constant $c\in\mathbb{N}$ and a finite
Turing machine that decides $L$ using at most $c\,s(n)$ space on inputs of length $n$.
\end{definition}

\begin{theorem}[Monotonicity of space classes]
\lean{Complexity.SPACE.mono}
\label{Complexity.SPACE.mono}
\leanok
\uses{Complexity.SPACE, Turing.FinTM.ComputesInSpace.mono}
\difficulty{4}
Let $s_1,s_2:\mathbb{N}\to\mathbb{N}$. If $s_1(n)\le s_2(n)$ for every $n\in\mathbb{N}$,
then every language decidable in space $s_1$ up to a constant factor is also decidable
in space $s_2$ up to a constant factor.
\end{theorem}

\begin{theorem}[Indexed language]
\lean{Complexity.indexLang}
\label{Complexity.indexLang}
\leanok
\uses{Turing.pairEncode}
\difficulty{1}
For a predicate $p$ on pairs consisting of a Boolean string $x$ and a natural number
$i$, the indexed language contains exactly those strings that encode a pair $(x,	ext{the
binary expansion of }i)$ for which $p(x,i)$ holds.
\end{theorem}

\begin{definition}[Monotonicity of the logarithmic bound]
\lean{Complexity.logSpace_mono}
\label{Complexity.logSpace_mono}
\leanok
\uses{Complexity.logSpace}
\difficulty{3}
If $a,b\in\mathbb{N}$ and $a\le b$, then $\ell(a)\le\ell(b)$, where
$\ell(n)=\lfloor\log_2 n\rfloor+1$ for positive $n$ and $\ell(0)=1$.
\end{definition}

%%% added by the blueprint pipeline
\subsection*{Space Complexity --- Foundations: Further results}

\begin{definition}[Computing within a space bound]
\lean{Turing.FinTM.ComputesInSpace}
\label{Turing.FinTM.ComputesInSpace}
\leanok
\uses{Turing.FinTM, Turing.FinTM.ComputesInTime, Turing.MultiTapeTM.initCfg, Turing.MultiTapeTM.spaceUsed}
A finite Turing machine $M$ computes a function $f$ within space bound $s$ if, for every
Boolean string $x$, there is a time $t$ at which $M$ outputs $f(x)$ and uses at most
$s(|x|)$ space, measured from its initial configuration on $x$.
\end{definition}

\begin{definition}[Deciding within a space bound]
\lean{Turing.FinTM.DecidesInSpace}
\statementsource{Sav70}{
  pdf-c8033d2fd705-p002-b002,
  pdf-c8033d2fd705-p003-b001
}
\label{Turing.FinTM.DecidesInSpace}
\leanok
\uses{Turing.FinTM, Turing.FinTM.ComputesInSpace, Turing.MultiTapeTM.indicator}
A finite Turing machine $M$ decides a language $L$ of Boolean strings within space bound
$s$ if, on every Boolean input string $x$, it computes the one-symbol Boolean indicator
of whether $x$ belongs to $L$, using at most $s(|x|)$ space.
\end{definition}

\begin{theorem}[Monotonicity of space bounds]
\lean{Turing.FinTM.ComputesInSpace.mono}
\label{Turing.FinTM.ComputesInSpace.mono}
\leanok
\uses{Turing.FinTM, Turing.FinTM.ComputesInSpace}
\difficulty{4}
Let $M$ be a finite Turing machine over the Boolean alphabet, let $f$ map Boolean words
to Boolean words, and let $s,s' : \mathbb{N} \to \mathbb{N}$. If for every input word
$x$ the machine computes $f(x)$ on input $x$ using space at most $s(|x|)$, and $s(n) \le
s'(n)$ for every $n$, then for every input word $x$ it computes $f(x)$ using space at
most $s'(|x|)$.
\end{theorem}

\begin{theorem}[Zero space with no work tapes]
\lean{Turing.FinTM.computesInSpace_zero_of_k_eq_zero}
\label{Turing.FinTM.computesInSpace_zero_of_k_eq_zero}
\leanok
\uses{Turing.FinTM, Turing.FinTM.ComputesInSpace, Turing.FinTM.ComputesInTime, Turing.MultiTapeTM.spaceUsed_zero_tapes_eq_zero}
\difficulty{4}
Let $M$ be a finite Turing machine over the Boolean alphabet and let $f$ map Boolean
words to Boolean words. If $M$ has no work tapes and for every Boolean input $x$ there
exists a time at which it halts with output $f(x)$, then $M$ computes $f$ using space
bounded by zero on every input.
\end{theorem}

\begin{theorem}[Monotonicity under constant scaling]
\lean{Complexity.SPACE.mono_const}
\label{Complexity.SPACE.mono_const}
\leanok
\uses{Complexity.SPACE, Turing.FinTM.ComputesInSpace.mono}
\difficulty{4}
Let $s_1,s_2:\mathbb{N}\to\mathbb{N}$ and let $a\in\mathbb{N}$. If $s_1(n)\le a s_2(n)$
for every $n\in\mathbb{N}$, then every language decidable in space $s_1$ up to a
constant factor is also decidable in space $s_2$ up to a constant factor.
\end{theorem}

\begin{definition}[Logarithmic space bound]
\lean{Complexity.logSpace}
\statementsource{Vad12}{
  pdf-bf2ce5422088-p294-b002,
  pdf-bf2ce5422088-p294-b003
}
\label{Complexity.logSpace}
\leanok
For each natural number $n$, define its logarithmic space bound to be
$\lfloor\log_2 n\rfloor+1$, where $\lfloor\log_2 n\rfloor$ denotes the integer part of
the base-$2$ logarithm, with value $0$ when $n=0$.
\end{definition}

\begin{definition}[Logarithmic-space languages]
\lean{Complexity.LOGSPACE}
\statementsource{Vad12}{
  pdf-bf2ce5422088-p037-b001,
  pdf-bf2ce5422088-p294-b002,
  pdf-bf2ce5422088-p294-b003
}
\label{Complexity.LOGSPACE}
\leanok
\uses{Complexity.SPACE, Complexity.logSpace}
The class of languages over the Boolean alphabet decidable by a Turing machine using
space at most logarithmic in the input length.
\end{definition}

\begin{definition}[Implicit logarithmic-space computability]
\lean{Complexity.ImplicitlyLogspaceComputable}
\label{Complexity.ImplicitlyLogspaceComputable}
\leanok
\uses{Complexity.LOGSPACE, Complexity.indexLang}
A function $f$ from Boolean strings to Boolean strings is implicitly logarithmic-space
computable if there are natural numbers $C,c$ such that $|f(x)|\le C(|x|+1)^c$ for every
input $x$, and both of the following languages are in logarithmic space: the language of
encodings of $(x,i)$ for which the $i$th output bit of $f(x)$ is true (taking an
out-of-range bit to be false), and the language of encodings of $(x,i)$ for which
$i<|f(x)|$.
\end{definition}
